\documentclass[11pt]{article}

\usepackage[margin=1in]{geometry}
\usepackage{amsmath,amssymb,amsthm}
\usepackage{mathtools}
\usepackage{hyperref}
\usepackage{xcolor}
\usepackage{enumitem}

\hypersetup{colorlinks=true, linkcolor=blue, citecolor=blue, urlcolor=blue}

\newcommand{\R}{\mathbb{R}}
\newcommand{\N}{\mathbb{N}}
\newcommand{\Rstar}{\mathbb{R}^{\star}}
\newcommand{\eps}{\varepsilon}
\let\straightepsilon\epsilon
\renewcommand{\epsilon}{\varepsilon}
\newcommand{\st}{\operatorname{st}}
\newcommand{\Var}{\operatorname{Var}}

\setlist[description]{leftmargin=0pt, style=unboxed, font=\normalfont\bfseries}

\title{Twenty Exercises in Probability on Ordinary Intervals: Solutions\\
\large Densities, $dx=\eps$, and $\omega$ as the Dirac delta}
\author{Exercises for engineers, in the hyperreal field $\Rstar$}
\date{\today}

\begin{document}
\maketitle


\section*{The two rules you need}

Nothing here is counted, and no hyperfinite sample space is built. The sample
space is an ordinary real interval such as $[0,1]$; all of the hyperreal
content sits in the integral. Fix a positive infinitesimal $\eps$, put
$\omega = 1/\eps$, and take the resolution of the line to be $\eps$: a point
$y$ \emph{is} the half-open dot $[y, y+\eps)$. The integral is then the
hyperfinite left-endpoint Riemann sum with $dx = \eps$,
\[
  \int_{[a,b)} f(x)\,dx
  := \sum_{k=0}^{(b-a)\omega-1} f(a+k\eps)\,\eps ,
\]
evaluated in $\Rstar$ and \emph{not} followed by $\st$. A probability is the
integral of a density, $P(E)=\int_E p$. Because a point is one dot, this
specializes to the identity used on nearly every page:
\[
  \boxed{\;P(\{y\}) = p(y)\,\eps\;}
\]
In particular a uniform law on $[0,1)$ has $p \equiv 1$, so hitting an exact
number has probability $\eps > 0$, not zero --- and $\int_{[0,1)} 1\,dx
= \omega\eps = 1$ exactly, so ordinary interval probabilities are untouched.

An atom of mass $a$ is not a second mechanism bolted onto the density. It is
the density \emph{value} $p(y)=a\omega$, since then $P(\{y\})=a\omega\eps=a$.
The Dirac delta is therefore an ordinary function taking the value $\omega$ on
one dot, and $\st$ is applied only where it is explicitly wanted.

\paragraph{The other model.} These exercises never count anything. The
companion curriculum answers the same questions by counting a hyperfinite
sample space $\{0,\dots,\omega-1\}$ instead, and its solutions are
machine-checked in Lean~4.
\begin{center}\small\url{https://files.pannous.com/hyperreal-exercises.pdf}\end{center}

\paragraph{Theory.} The axioms, the model, and the underlying calculus are
\emph{not} repeated here. They are developed in the companion paper,
\href{https://files.pannous.com/hyperreals-light.pdf}{\emph{Hyperreal Numbers: $\eps\cdot\omega=1$ --- An
Algebraic, Computable Introduction to Infinitesimal Calculus}}
(\url{https://files.pannous.com/hyperreals-light.pdf}), with the machine-checked Lean~4 development at
\url{https://github.com/pannous/hyper-lean}.

\paragraph{Readiness labels.} Each exercise is tagged by what the present
Lean~4 formalization already supports.
\begin{description}
\item[Now] the value and its derivation use only single-dot integrals and
$\Rstar$ arithmetic, which the concrete representation already has.
\item[Partial] the closed form is exact and short, but the derivation sums over
the $(b-a)\omega$ dots of an interval, for which there is no index type yet.
\item[Research] additionally needs transcendental densities, a proved
$\st$-compatibility theorem, or a genuine conditional-law construction.
\end{description}

\paragraph{Exercises.} The problems are stated in full in the companion
booklet \href{https://files.pannous.com/hyperreal-integral-exercises.pdf}{\emph{Twenty Exercises in Probability on Ordinary
Intervals}}:
\begin{center}\small\url{https://files.pannous.com/hyperreal-integral-exercises.pdf}\end{center}
Each problem is restated below, so this booklet is self-contained.

\paragraph{No checked kernels yet.} Unlike the counting curriculum, this
framework has \emph{no} Lean~4 formalization: \texttt{Hyper/HyperList.lean}
offers no integral, no density type, and above all no index type of size
$\omega$ to sum over. The solutions below are worked algebra, not machine-checked
theorems, and the readiness labels measure distance from the current code
rather than proofs already obtained. Anything needing a genuine
$\sum_{k<\omega}$ is labelled \textbf{Research} however short its algebra is.


\section*{Solutions}


\subsection*{Exercise 1 --- The probability of hitting an exact number \hfill \normalfont\small[Now]}

\textbf{Problem.} Let $X$ be uniform on $[0,1)$, so $p(x) = 1$. Compute $P(X = y)$
for a specified $y$, prove that it is positive, and give its standard part.

\textbf{Solution.} The point $y$ is one dot of width $\eps$, so the integral has
a single term:

\[
P(X=y)=\int_{[y,y+\epsilon)} 1\,dx=1\cdot\epsilon=\epsilon>0,
\qquad \operatorname{st}(\epsilon)=0 .
\]

No limit and no counting argument is involved: $\eps$ appears because
$dx = \eps$.

\textbf{Ingredients and readiness.} The integral axiom on one dot, positivity of
$\eps$, and $\st(\eps) = 0$. \textbf{Now.}


\subsection*{Exercise 2 --- Finitely many targets \hfill \normalfont\small[Now]}

\textbf{Problem.} With $X$ uniform on $[0,1)$, let $A$ be a set of $r$ specified
points, $r$ a standard positive natural number. Find $P(A)$ and the ratio
$P(A)/P(\{y\})$.

\textbf{Solution.} The dots are disjoint, so the integral is additive over them:

\[
P(A)=r\epsilon,\qquad \frac{P(A)}{P(\{y\})}=r .
\]

\textbf{Ingredients and readiness.} Finite additivity over disjoint dots and
cancellation of $\eps$. \textbf{Now.}


\subsection*{Exercise 3 --- Intervals are unharmed, and the total is exactly one \hfill \normalfont\small[Now]}

\textbf{Problem.} For $X$ uniform on $[0,1)$ and standard $0 \le  c < d \le  1$,
compute $P([c,d))$ and verify $P([0,1)) = 1$ exactly.

\textbf{Solution.} The interval $[c,d)$ consists of $(d-c)\cdot \omega$ dots, each
contributing $\eps$:

\[
P([c,d))=(d-c)\omega\cdot\epsilon=d-c ,
\]

and with $c = 0$, $d = 1$ this is $\omega\cdot \eps = 1$. The gauging axiom is
exactly what makes the total mass come out without an error term.

\textbf{Ingredients and readiness.} A hyperfinite sum of a constant over the dots
of an interval, and $\eps\cdot \omega = 1$. \textbf{Now:} checked as
$integral (constant 1) 0 1 = 1$ and $prob (1/4) (3/4) = 1/2$.


\subsection*{Exercise 4 --- Half-open versus closed, and what a normalization costs \hfill \normalfont\small[Partial]}

\textbf{Problem.} Integrate the constant $1$ over the \emph{closed} interval $[0,1]$. Then
find the density of the uniform law on $[0,1]$ and recompute $P(\{0\})$.

\textbf{Solution.} The closed interval contains one more dot than the half-open
one, so

\[
\int_{[0,1]} 1\,dx=(\omega+1)\epsilon=1+\epsilon\neq 1 .
\]

A genuine density on $[0,1]$ must therefore be renormalized,

\[
p=\frac1{1+\epsilon}=1-\epsilon+\epsilon^{2}-\cdots,
\qquad
P(\{0\})=p\,\epsilon=\epsilon-\epsilon^{2}+\cdots .
\]

Both conventions are consistent; only the half-open one gives $p = 1$.

\textbf{Ingredients and readiness.} Dot counting of an interval, inversion of
$1 + \eps$, and series expansion. \textbf{Partial;} the inversion of a
multi-term denominator is the same gap recorded in the counting curriculum.


\subsection*{Exercise 5 --- Conditioning on an ordinary interval \hfill \normalfont\small[Now]}

\textbf{Problem.} With $X$ uniform on $[0,1)$ and $y$ in $[a,b)\subset [0,1)$,
compute $P(X = y\mid a \le  X < b)$.

\textbf{Solution.} \[
P(X=y \mid a\le X<b)=\frac{\epsilon}{b-a},
\]

which is the point probability of the uniform law on $[a,b)$, as conditioning
should give.

\textbf{Ingredients and readiness.} Conditional probability as field division and
division of $\eps$ by a nonzero standard number. \textbf{Now.}


\subsection*{Exercise 6 --- A nonuniform density is read off at the point \hfill \normalfont\small[Now]}

\textbf{Problem.} Let $X$ have the triangular density $p(x) = 2x$ on $[0,1)$.
Verify the normalization, compute $P(X = y)$, and compare two points $y1$ and
$y2$.

\textbf{Solution.} Under the canonical midpoint rule the normalization is exact,
$integral of 2x over [0,1) = 1$; under the left rule it would be $1 - \eps$,
requiring $p(x) = 2x/(1-\eps)$ instead. In either case

\[
P(X=y)=p(y)\,\epsilon,
\qquad
\frac{P(X=y_1)}{P(X=y_2)}=\frac{p(y_1)}{p(y_2)}=\frac{y_1}{y_2}.
\]

Every point still has infinitesimal probability, but their ratios are ordinary
real numbers.

\textbf{Ingredients and readiness.} Single-dot integration and division of
infinitesimals of equal order. \textbf{Now:} under the canonical midpoint rule the
triangular density is exactly normalized, and both the normalization and
$P(\{y\}) = 2y\cdot \eps$ are checked.


\subsection*{Exercise 7 --- An atom is a density value, not a second mechanism \hfill \normalfont\small[Now]}

\textbf{Problem.} Let $p$ be the density that equals $\omega/2$ on the dot at $0$
and $1/2$ on $(0,1)$. Show that $p$ is a probability density, compute
$P(\{0\})$ and $P(\{y\})$ for $y \neq  0$, and compare their orders.

\textbf{Solution.} The total mass splits into the atom's single dot and the rest:

\[
\int p=\frac\omega2\cdot\epsilon+\frac12\cdot 1=\frac12+\frac12=1 .
\]

Then

\[
P(\{0\})=\frac\omega2\cdot\epsilon=\frac12,
\qquad
P(\{y\})=\frac\epsilon2 \quad (y\neq0),
\qquad
\frac{P(\{0\})}{P(\{y\})}=\omega .
\]

The atom is infinitely more likely than an ordinary point, which is the exact
algebraic content of the word ``atom''.

\textbf{Ingredients and readiness.} An $\omega$-valued density, splitting an
integral over disjoint regions, and order comparison. \textbf{Now:} the mixed law is
checked end to end, including that its total mass is exactly $1$.

The atom here \emph{is} the Dirac delta of Exercise 8, scaled: $a\cdot \delta$. One
object, one mechanism --- that is what $\omega$ unifies. What is two-celled is not
$\delta$ but the central difference \emph{operator} of Exercise 8, whose $2\cdot \eps$
stencil returns $\delta$ smeared over the halo.


\subsection*{Exercise 8 --- The Dirac delta is the derivative of the step, exactly \hfill \normalfont\small[Now]}

\textbf{Problem.} Let $H$ be the Heaviside step, $H(x) = 1$ for $x \ge  0$ and $0$
otherwise. Compute the algebraic derivative $dH(x) = (H(x+\eps)-H(x))/\eps$
at every $x$, and integrate the result over $[-1,1)$.

\textbf{Solution.} The difference quotient vanishes unless the dot straddles the
jump, and

\[
\partial H(x)=\omega \quad\text{for } x\in[-\epsilon,0),
\qquad \partial H(x)=0 \text{ otherwise},
\]

so $dH = \delta_{-\eps}$ and

\[
\int_{[-1,1)}\partial H(x)\,dx=\omega\cdot\epsilon=1 .
\]

The fundamental theorem of calculus therefore holds for the step function on
the nose, with no distributions and no limits.

\textbf{Ingredients and readiness.} The algebraic derivative, a one-dot integral,
and $\eps\cdot \omega = 1$. \textbf{Now.}


\subsection*{Exercise 9 --- Three orders of point probability \hfill \normalfont\small[Now]}

\textbf{Problem.} At a point $y$, consider densities $p(y) = a\cdot \omega$, $p(y) = b$,
and $p(y) = c\cdot \eps$ with standard positive $a,b,c$. Compute $P(\{y\})$ in
each case and order the three results.

\textbf{Solution.} By $P(\{y\}) = p(y)\cdot \eps$,

\[
a\omega\cdot\epsilon=a,
\qquad
b\cdot\epsilon=b\epsilon,
\qquad
c\epsilon\cdot\epsilon=c\epsilon^{2},
\]

so $a > b\cdot \eps > c\cdot \eps^2$ regardless of the standard coefficients. A
point's rarity order is exactly the negative of the density's order.

\textbf{Ingredients and readiness.} Multiplication by $\eps$, and comparison by
leading exponent. \textbf{Now.}


\subsection*{Exercise 10 --- The dart, with no grid at all \hfill \normalfont\small[Partial/Now]}

\textbf{Problem.} Let $(X,Y)$ be uniform on the unit square, $p = 1$. Compute the
probability of one specified point.

\textbf{Solution.} The area element is $dA = dx\cdot dy = \eps^2$, so a single dot of
the plane has

\[
P(\{(x_0,y_0)\})=1\cdot\epsilon^{2}=\epsilon^{2}.
\]

\textbf{Ingredients and readiness.} A product integral and $\eps^2$. \textbf{Partial}
for the product construction; the value itself is \textbf{Now}.


\subsection*{Exercise 11 --- A segment beats a point by exactly $\omega$ \hfill \normalfont\small[Partial]}

\textbf{Problem.} In the same square, compute $P(Y = X)$, the probability that the
two coordinates agree, and compare it with Exercise 10.

\textbf{Solution.} Condition on $X$ and integrate the one-dimensional point
probability:

\[
P(Y=X)=\int_0^1 P(Y=x)\,dx=\int_0^1 \epsilon\,dx=\epsilon,
\qquad
\frac{P(Y=X)}{P(\{(x_0,y_0)\})}=\omega .
\]

Each integration that is \emph{not} performed costs one factor of $\eps$; this
is the codimension law of the counting model, derived instead of postulated.

\textbf{Ingredients and readiness.} Iterated integration and monomial division.
\textbf{Partial.}


\subsection*{Exercise 12 --- A uniform law on the whole line \hfill \normalfont\small[Partial]}

\textbf{Problem.} Take the ambient line to be $[-\omega, \omega)$. Find the uniform
density on it, compute $P(\{y\})$ and $P([a,b))$, and compare $P(\{y\})$ with the
point probability of the uniform law on $[0,1)$.

\textbf{Solution.} The ambient length is $2\cdot \omega$, so normalization forces
$p = 1/(2\cdot \omega) = \eps/2$ and

\[
P(\{y\})=\frac\epsilon2\cdot\epsilon=\frac{\epsilon^{2}}2,
\qquad
P([a,b))=(b-a)\frac\epsilon2 .
\]

Comparing with Exercise 1,

\[
\frac{P_{\mathbb R}(\{y\})}{P_{[0,1)}(\{y\})}=\frac\epsilon2 ,
\]

infinitesimal: a point is infinitely rarer when the target is the whole line.
Any finite interval likewise gets only infinitesimal probability, which is the
correct and previously unavailable statement.

\textbf{Ingredients and readiness.} The ambient-line convention, inversion of
$2\cdot \omega$, and order comparison. \textbf{Partial;} the choice of $[-\omega, \omega)$
over the one-sided convention must be stated, not derived.


\subsection*{Exercise 13 --- Exact moments of the uniform law \hfill \normalfont\small[Now]}

\textbf{Problem.} For $X$ uniform on $[0,1)$, compute $E[X]$, $E[X^2]$, and
$\Var(X)$ from the integral, without taking any limit.

\textbf{Solution.} The finite power-sum identities give exact values, and the
sampling convention is visible in them. Under the left rule,

\[
E[X]=\sum_{k<\omega}k\epsilon\cdot\epsilon
=\epsilon^{2}\frac{\omega(\omega-1)}2=\frac12-\frac\epsilon2,
\qquad
E[X^{2}]=\frac13-\frac\epsilon2+\frac{\epsilon^{2}}6 ,
\]

while under the canonical midpoint rule the first-order bias disappears:

\[
E[X]=\frac12 \ \text{exactly},
\qquad
E[X^{2}]=\frac13-\frac{\epsilon^{2}}{12},
\qquad
\operatorname{Var}(X)=\frac1{12}-\frac{\epsilon^{2}}{12}.
\]

Taking $st$ recovers $1/2$ and $1/12$ only after the exact corrections have
been displayed.

\textbf{Ingredients and readiness.} Hyperfinite power sums and substitution
$\omega = 1/\eps$. \textbf{Now:} the power-sum recursion evaluates these exactly
without an index type; the left-rule values above and the midpoint values
$E[X] = 1/2$, $E[X^2] = 1/3 - \eps^2/12$ are both checked.


\subsection*{Exercise 14 --- The integral convention is visible at order $\eps$ \hfill \normalfont\small[Now/Research]}

\textbf{Problem.} For a standard $f$ on $[0,1)$, compare the left-endpoint,
right-endpoint, and midpoint sums. Apply the result to $f(x) = x$.

\textbf{Solution.} The right and left sums differ only in their end terms,

\[
\int^{\text{right}}f-\int^{\text{left}}f=\epsilon\,(f(1)-f(0)),
\]

and for $f(x) = x$ the midpoint rule gives exactly $1/2$, against
$1/2 - \eps/2$ and $1/2 + \eps/2$. All three agree after $st$.

Conclusion: ``the integral'' is not one object in this framework. A convention
must be fixed, and every exercise above uses the left-endpoint one.

\textbf{Ingredients and readiness.} Telescoping of a hyperfinite sum and
$st$-invariance. \textbf{Now} for the $f(x) = x$ instance: all three rules and their
exact difference $\eps$ are checked. \textbf{Research} for the general
statement.


\subsection*{Exercise 15 --- A mixed law needs no case split \hfill \normalfont\small[Research]}

\textbf{Problem.} With the density of Exercise 7 --- mass $1/2$ at the point $0$ and
$1/2$ spread uniformly --- compute $E[X]$ and the distribution function
$F(t) = P(X < t)$.

\textbf{Solution.} The atom contributes at $x = 0$ and the continuous part is
Exercise 13 halved:

\[
E[X]=0\cdot\frac12+\frac12\left(\frac12-\frac\epsilon2\right)
=\frac14-\frac\epsilon4 ,
\]

\[
F(t)=\frac12+\frac t2 \quad\text{for } 0<t\le1,
\qquad F(0)=0,\quad F(\epsilon)=\frac12 .
\]

The jump of $F$ at $0$ is the atom's mass, produced by integrating an
$\omega$-valued density over one dot rather than by adding a point weight.

\textbf{Ingredients and readiness.} Integration of a density with an $\omega$ value,
and the moment computation of Exercise 13. \textbf{Research,} for the same reason
as Exercise 13.


\subsection*{Exercise 16 --- Point probability is not invariant, and should not be \hfill \normalfont\small[Now/Partial]}

\textbf{Problem.} Let $X$ be uniform on $[0,1)$ and $Y = 2X$. Find the density of
$Y$, compute $P(Y = y)$, and reconcile it with $P(X = y/2) = \eps$.

\textbf{Solution.} $Y$ is uniform on $[0,2)$, so $p_Y = 1/2$ and

\[
P(Y=y)=\frac12\epsilon .
\]

There is no contradiction: the event $\{Y = y\}$ is $\{X = y/2\}$ only up to
resolution, since the map doubles dot widths, so one dot of $Y$ corresponds to
half a dot of $X$. In general $p_Y(y) = p_X(x)/|g'(x)|$ for $y = g(x)$, with
the algebraic derivative, and $P(\{y\}) = p_Y(y)\cdot \eps$ follows.

\textbf{Ingredients and readiness.} Change of variables with the algebraic
derivative, and the observation that dots are not preserved by non-isometries.
\textbf{Now} for the linear case; \textbf{Partial} in general.


\subsection*{Exercise 17 --- An exponential law \hfill \normalfont\small[Research]}

\textbf{Problem.} Let $p(x) = \lambda\cdot e^{-\lambda\cdot x}$ on $[0, \omega)$ with standard
$\lambda > 0$. Compute $P(X = y)$ and $P(X \ge  t)$ for standard $t$, and discuss
the normalization.

\textbf{Solution.} Pointwise the rule is immediate,

\[
P(X=y)=\lambda e^{-\lambda y}\epsilon ,
\]

so a point far out in the tail is still infinitesimal but by an ordinary
factor $e^{-\lambda\cdot y}$ less likely than one near $0$. The tail and the total
are hyperfinite geometric sums, and both carry a discretization correction:

\[
\int_{[0,\omega)} p
=\frac{\lambda\epsilon}{1-e^{-\lambda\epsilon}}\left(1-e^{-\lambda\omega}\right)
=1+\frac{\lambda\epsilon}2+O(\epsilon^{2}),
\qquad
P(X\ge t)=e^{-\lambda t}\left(1+\frac{\lambda\epsilon}2+O(\epsilon^{2})\right).
\]

So the left-endpoint convention \emph{over}-counts by $\lambda\cdot \eps/2$, and the
density must be renormalized by that factor to be exact; the truncation at
$\omega$ costs only $e^{-\lambda\cdot \omega}$, an infinitesimal smaller than every
power of $\eps$. The classical answers $1$ and $e^{-\lambda\cdot t}$ are the
standard parts.

\textbf{Ingredients and readiness.} A transcendental density, hyperfinite geometric
summation, and comparison of $e^{-\lambda\cdot \omega}$ with all powers of $\eps$.
\textbf{Research:} \texttt{Hyper/HyperTranscendental.lean} has $hexp$ but it is not
connected to any integral.


\subsection*{Exercise 18 --- Bayes with an atom and a density in one formula \hfill \normalfont\small[Research]}

\textbf{Problem.} A parameter $T$ is $0$ with probability $1/2$ and otherwise
uniform on $(0,1)$. Given $T = t$, an observation $Z$ has a standard density
$f(z | t)$. Compute the posterior $P(T = 0\mid Z = z)$.

\textbf{Solution.} Write the prior as one density: $pi = (\omega/2)$ on the dot at
$0$ and $1/2$ on $(0,1)$. Every term of Bayes' rule then carries one factor
$\eps$ from the observed event $\{Z = z\}$, which cancels, and the atom's
$\omega$ cancels the $\eps$ of its own dot, leaving

\[
P(T=0\mid Z=z)
=\frac{\tfrac12 f(z\mid 0)}
       {\tfrac12 f(z\mid 0)+\tfrac12\int_0^1 f(z\mid t)\,dt } .
\]

The answer is finite and ordinary --- the infinitesimals cancel --- but it was
obtained without ever conditioning on a classically null event by fiat.

\textbf{Ingredients and readiness.} Joint densities, cancellation of $\eps$
between numerator and denominator, and inversion of a finite value.
\textbf{Research:} needs a joint-density construction and the $st$-compatibility of
the inner integral.


\subsection*{Exercise 19 --- The Borel--Kolmogorov paradox dissolves \hfill \normalfont\small[Research/Partial]}

\textbf{Problem.} For $(X,Y)$ uniform on the unit square, condition on the diagonal
$D = \{Y = X\}$. Compute $P(X\in [a,b)\mid D)$ and explain why the classical
paradox does not arise.

\textbf{Solution.} By Exercise 11, $P(D) = \eps$. The part of the diagonal over
$[a,b)$ has probability $(b-a)\cdot \eps$, so

\[
P\big(X\in[a,b)\mid D\big)=\frac{(b-a)\epsilon}{\epsilon}=b-a .
\]

The conditional law is uniform, and there was no choice to make: the diagonal
is an ordinary event of positive probability, so $P(·|D)$ is the ordinary
quotient. A different description of the same set of dots gives the same
value, because it is the same event. The paradox was an artifact of dividing
by zero.

\textbf{Caution:} The dissolution is real but narrow: it holds for the dot resolution fixed
in the foundations note. A \emph{different} resolution convention describes a
different event, and the classical ambiguity reappears as the choice of
convention --- now explicit and stated, instead of hidden.

\textbf{Ingredients and readiness.} Conditioning on an infinitesimal-probability
event, and the diagonal integral of Exercise 11. \textbf{Research} as a general
theorem; the displayed algebra is \textbf{Partial}.


\subsection*{Exercise 20 --- What must be proved before this is a measure theory \hfill \normalfont\small[Research]}

\textbf{Problem.} Identify precisely what is needed to prove: (a) that the standard
part of the hyperfinite integral is the classical Riemann integral, for every
standard Riemann-integrable $f$; (b) that $P$ is finitely additive over arbitrary disjoint unions of dots;
(c) that $P$ is \emph{not} countably additive in the classical sense, and why that
is not a defect.

\textbf{Ingredients and readiness.} Transfer or explicit estimates, a theory of
dot-unions, and a Loeb-style comparison with classical measure.
\textbf{Research:} this is the capstone, and none of its prerequisites exist in the
repository today.


\section*{What to build first}

The exercises fall into three implementation stages. Stage one is a
single-dot integral and the identity $P(\{y\})=p(y)\eps$, which needs nothing
beyond the arithmetic already proved for $\Rstar$. Stage two is an index type
ranging over the $(b-a)\omega$ dots of an interval, with the constant and
power sums; that single addition turns most \textbf{Partial} labels into
\textbf{Now}. Stage three is the analytic layer --- transcendental densities,
$\st$-compatibility with the classical Riemann integral, and conditional laws
on infinitesimal-probability events --- which is where the remaining
\textbf{Research} exercises live.

The conventions used throughout (half-open dots, left endpoints,
$\int\delta = 1$ on one dot, the ambient line $[-\omega,\omega)$) are choices,
not consequences. \texttt{notes/integral/integral-probability-foundations.md}
lists each one together with what the alternative would change; every
alternative changes answers at order $\eps$, which is exactly the order this
framework exists to discuss.


\end{document}
